\documentclass[10pt,a4paper]{amsart}
\usepackage[margin=19mm]{geometry}
\usepackage{amsmath,amssymb,mathtools}
\usepackage[scr=rsfs,cal=euler]{mathalfa}
\usepackage{xfrac}
\usepackage[unicode,hidelinks,bookmarks=false]{hyperref}
\newcommand{\ie}{i.e.\ }
\newcommand{\cf}{cf.\ }
\newcommand{\resp}{respectively\ }
\newcommand{\Subsection}{\subsection}
\providecommand{\nobreakdash}{\nobreak\hbox{-}}
\setlength{\emergencystretch}{2em}
\multlinegap0pt
\usepackage{amssymb}
\usepackage{enumerate}
\usepackage{tikz}
\newtheorem{lemma}{Lemma}
\newtheorem{proposition}{Proposition}
\usepackage{cleveref}
\crefname{lemma}{Lemma}{Lemmas}
\crefname{proposition}{Proposition}{Propositions}
\newcommand{\ZZ}{{\mathbb Z}}
\DeclareMathOperator{\lcm}{lcm}
\title{Standard Veronese degree and star configurations}
\author{Th\'ai Th\`anh Nguy$\tilde{\text{\^E}}$n}
\date{Source: arXiv:2607.03500v1 (CC BY 4.0)}
\begin{document}
\maketitle

\noindent\emph{Source and licence.} Standard Veronese degree and star configurations. Version arXiv:2607.03500v1 (CC BY 4.0), distributed by arXiv under the Creative Commons Attribution 4.0 International licence, http://creativecommons.org/licenses/by/4.0/, as recorded in the arXiv licence metadata for this exact version and shown on the abstract page https://arxiv.org/abs/2607.03500v1. Full licence notice: \texttt{licenses/arxiv-2607.03500-CC-BY-4.0.txt}.\par\smallskip

\noindent\textit{Editorial scope.} Counted unit: the article's proposition prop.minSol (source0.tex lines 304--307). The article's complete original proof of it is delivered with it (source0.tex lines 308--335, one \texttt{proof} environment of 28 lines). No mathematics is written, repaired or re-derived: the statement and the proof are the article's own lines, in the article's own notation, and no retained line is substituted or re-flowed.  Retained uncounted as the source-local context the proof needs: lines 265--268; lines 280--283; lines 292--295.  Nothing was omitted from any retained span, so the package carries no omission record.  No retained line contains a citation, so the wrapper carries no bibliography and no bibliography entry.  No retained line contains a figure environment, an image-inclusion command or drawing code, so no image is delivered and the package declares no asset dependency.  Editorial formatting only, in addition: an AMS \texttt{amsart} preamble that restates the article's own declarations and macros used by the retained lines, namely the environments \texttt{lemma}, \texttt{proposition} on separate counters, together with the standard packages those lines need.  The retained environments print in this document as Lemma 1, Proposition 1 in order of appearance, whereas the article prints the same items with the numbering of its own full document.  The complete native source of the article as distributed in the arXiv e-print for this exact version is bundled unchanged as \texttt{sources/204397/source0.tex}.
\begin{lemma}
    \label{lem.xi<n}
    Let $k_1,k_2,\ldots,k_s \in \ZZ_{>0}$ and $d=\lcm(k_1,\ldots ,k_s)$. Suppose that $(a_1,a_2,\ldots ,a_s, b)$ with $b\ge 2$ is a minimal solution of the equation $k_1x_1+k_2x_2+\ldots +k_sx_s = dy$ (in variables $x_1,\ldots,x_s,y$). Then $a_i\le \dfrac{d}{k_i}-1$ for all $i=1,\ldots ,s$. 
\end{lemma}

\begin{lemma}
    \label{lem.combIneq}
    Let $d_n=\lcm(1,2,\ldots , n)$. Then $d_n >\dfrac{(n-2)(n+1)n}{4}$ for all $n\ge 1$.
\end{lemma}

\begin{lemma}
    \label{pigeonhole}
    Let $a_1,a_2,\ldots ,a_n$ be a sequence of integers. Then for any $1\le m \le n$, there exists a subsequence $a_{i_1},a_{i_2},\ldots ,a_{i_k}$ with $k\ge n-m+1$ elements such that $m \mid (a_{i_1}+a_{i_2}+\ldots +a_{i_k})$.
\end{lemma}

\begin{proposition}
    \label{prop.minSol}
    For each $n\ge 1$, the equation $x_1+2x_2+\ldots +nx_n=d_ny$ where $d_n=\lcm(1,2,\ldots , n)$ satisfies condition $\mathcal{M}$.
\end{proposition}

\begin{proof}
    The case where $n=1$ is trivial. Assume that $n\ge 2$. Suppose that $(a_1,a_2,\ldots ,a_n,b)$ is a minimal solution of the equation $x_1+2x_2+\ldots +nx_n=d_ny$ with $b\ge 2$. By \Cref{lem.combIneq}, $d_nb \ge 2d_n > (n-2)(1+2+\ldots +n)$, thus, there exists $i$ such that $a_i\ge n-1$. Moreover, by \Cref{lem.xi<n}, we have $a_i \le \dfrac{d_n}{i}-1$. It follows that 
    \[
    a_1+2a_2 +\ldots + \widehat{a_i} + \ldots +na_n \ge 2d_n - ia_i \ge d_n +i,
    \]
    where $\widehat{a_i}$ indicates the absence of $a_i$ in the sum. Let $S = \sum_{j \ne i} j a_j \ge d_n + i$. We construct a finite, strictly decreasing sequence of integers $S_0, S_1, \ldots, S_N$ as follows. Set $S_0 = S$. To obtain $S_{k+1}$ from $S_k$, we choose any index $j \in \{1, \ldots, n\} \setminus \{i\}$ for which the current coefficient of $j$ is strictly positive, and we reduce that coefficient by $1$. This process terminates in $N = \sum_{j \ne i} a_j$ steps when all coefficients reach zero, yielding $S_N = 0$. Since the sequence $S_0, S_1, \ldots, S_N$ must cross the value $d_n + i$ exactly once, along this specific sequence of reductions, there is a unique index $m \ge 1$ where $S_m < d_n + i$ while $S_{m-1} \ge d_n + i$.

    Let $a_1', \dots, \widehat{a_i'}, \dots, a_n'$ be the specific coefficients that correspond to the sum $S_m$. By construction, $0 \le a_j' \le a_j$ for all $j \ne i$, and moreover $S_{m-1} - S_m \le n$. Thus, we obtain
    \[
    d_n+i-n \le a_1'+2a_2'+\ldots + na_n' < d_n + i.
    \]

    Note that there is no $a_i'$ in the above sum. Let $a=a_1'+a_2'+\ldots +a_n'$. If $a\le i-1$, then
    \[
    d_n+i-n \le a_1'+2a_2'+\ldots + na_n'\le n(a_1'+a_2'+\ldots +a_n')\le n(i-1).
    \]
    But this implies that
    \[
    d_n \le i(n-1) \le ia_i \le i\left(\dfrac{d_n}{i}-1 \right) = d_n-i,
    \]
    a contradiction. Therefore, $a\ge i$. Consider the following sequence of a numbers: the number $1$ appearing $a_1'$ times, the number $2$ appearing $a_2'$ times, $\ldots$, the number $n$ appearing $a_n'$ times, and the number $i$ not appearing at all. Since $a\ge i$, by \Cref{pigeonhole}, there exists a subsequence of $a_1''$ numbers $1$, $a_2''$ numbers $2$, $\ldots$, and $a_n''$ numbers $n$ with $0\le a_j''\le a_j'$ for all $1\le j \le n, j\ne i$, $a_1''+a_2''+\ldots+ a_n'' \ge a-i+1$ such that $i \mid (a_1''+2a_2''+ \ldots+ na_n'')$. Furthermore, since the total numbers that are not in the subsequence is $a-(a_1''+a_2''+ \ldots+ a_n'') \le i-1$, it follows that
    \[
    a_1''+2a_2''+ \ldots+ na_n'' \ge d_n+i-n - (i-1)n =  d_n-i(n-1). 
    \]
    Note that we also have $a_1''+2a_2''+ \ldots+ na_n'' < d_n+i$. Let $a_i''=\dfrac{d_n-(a_1''+2a_2''+ \ldots+ na_n'')}{i}$, then $a_i'' \in \ZZ$ because $d_n$ is the $\lcm(1,2,\ldots,n)$, and $0 \le a_i'' \le n-1$. Thus, we can write
    $$(a_1,\ldots ,a_n,b) = (a_1'',\ldots, a_i'', \ldots,a_n'',1)+(a_1-a_1'',\ldots, a_i-a_i'', \ldots, a_n-a_n'',b-1),$$
    where each of the right hand side is a nonnegative solution to the equation $x_1+2x_2+\ldots +nx_n=d_ny$. This contradicts the assumption that $(a_1,a_2,\ldots ,a_n,b)$ is a minimal solution, and hence, finishes the proof.    
\end{proof}

\end{document}
